\section{Conclusion}
\label{sec:conclusion}

This study establishes a unified, comparative planetary magnetobiology framework for deep-space exploration, quantitatively mapping the remnant magnetic environments across 36 landing sites on the Moon and Mars and evaluating their biological implications. Our analysis demonstrates that all 13 candidate Artemis III South Polar landing sites on the Moon ($< \SI{1.0}{\nano\tesla}$) and ice-rich human habitat candidate sites in the Martian northern plains ($< \SI{20}{\nano\tesla}$) represent severe, near-null hypomagnetic field environments ($>2,500$-fold to $>60,000$-fold weaker than Earth's geomagnetic baseline). In contrast, Martian southern highland anomalies form localized mini-magnetospheres exceeding \SI{1500}{\nano\tesla} at ground level. Synthesizing empirical magnetobiology literature indicates that unmitigated hypomagnetic exposure poses distinct developmental hazards to crop plants, alters closed-loop habitat microbiomes, and synergistically accelerates musculoskeletal bone loss and cellular DNA repair perturbations in astronauts. We conclude that planetary magnetic field architecture is an essential environmental factor in space exploration, and we recommend integrating active Helmholtz coil magnetic restoration into extraterrestrial growth modules alongside standardized surface magnetometry on precursor landers.
